% 第3章 3.4.4：l^2 上 A(x_n)=(x_n/n) 的 Picard 条件。奇异值 sigma_n=1/n，u_n=v_n=e_n。
% g^(1)=(1/n^2)：系数 1/n^2，比值 |<u_n,g>|/sigma_n=1/n，平方可和，g^(1)=A(1/n)。
% g^(2)=(1/n)：系数 1/n 平方可和，但比值恒为 1，不平方可和，g^(2) 不在像中。n=1..15。
\begin{tikzpicture}[font=\figurefont\footnotesize]
\pgfplotsset{every axis/.append style={width=54mm,height=42mm,xmin=0.5,xmax=15.5,ymin=0,ymax=1.15,
  axis x line=bottom,axis y line=left,xtick={1,5,10,15},ytick={0,0.5,1},
  tick label style={font=\figurefont\scriptsize,text=FigureInk},
  xlabel={$n$},xlabel style={font=\figurefont\scriptsize},
  ymajorgrids,grid style={draw=FigureGrid,line width=.4pt},
  legend style={at={(0.97,0.97)},anchor=north east,font=\figurefont\scriptsize,cells={anchor=west}},
  title style={font=\figurefont\footnotesize}}}
\begin{axis}[name=A,title={(a) 数据系数$|\langle u_n,g\rangle|$}]
\addplot[FigureBlue,line width=1.1pt,mark=*,mark size=1pt,samples at={1,...,15}] {1/x^2};
\addplot[FigureYellow,line width=1.1pt,mark=square*,mark size=.9pt,samples at={1,...,15}] {1/x};
\legend{$g^{(1)}_n=1/n^2$,$g^{(2)}_n=1/n$}
\end{axis}
\begin{axis}[at={(A.east)},anchor=west,xshift=14mm,legend style={at={(0.97,0.62)},anchor=east},title={(b) 形式解系数$|\langle u_n,g\rangle|/\sigma_n$}]
\addplot[FigureBlue,line width=1.1pt,mark=*,mark size=1pt,samples at={1,...,15}] {1/x};
\addplot[FigureYellow,line width=1.1pt,mark=square*,mark size=.9pt,samples at={1,...,15}] {1};
\legend{$1/n$：平方可和,$1$：不平方可和}
\end{axis}
\end{tikzpicture}
